\headline{{\textbf{\Large\sffamily Annual General Meeting Summary:}}\
{\textbf{\large\sffamily November 2025}}}
\label{2025-11-agm-summary}

\noindent
Our Annual General Meeting brought together a strong turnout of members, both in person and online. Twenty-six attended at Whitton Community Centre, with a further two joining through the Discord chat. Despite a little echo on the audio, the hybrid set-up once again allowed broader participation, and the general sentiment across the membership was one of continued satisfaction. As reflected in conversations throughout the year, members value the club's friendly atmosphere, varied viewpoints, and steady organisational development. The steering committee’s work on activities, shows, and day-to-day administration was met with appreciation, and no concerns were raised about the direction of the club. It is heartening to note that members unable to attend also provided positive feedback, reinforcing the sense that the club is functioning well and delivering what people hope for: a welcoming, active community that supports bonsai growth, learning, and enjoyment.

\medskip
\topic{Committee and Governance}
\noindent
No new nominations were received this year, and the existing committee agreed to continue in their roles. A key proposal emerged from Tony: that the three major positions---Chairperson, Treasurer, and Secretary---commit to three-year terms in order to ensure stable leadership and consistent progress. This was unanimously approved, with Tony Ulatowski, Brendan Gibbs, and Nick Lloyd agreeing to serve these extended terms.
%
The remainder of the committee remains unchanged, covering vice-chair duties, events, multi-media, merchandise, and the newsletter, along with the various support roles that keep the club functioning smoothly. Notably, Mark E offered to assist with adding content to the Facebook page, strengthening the digital presence that many members rely on for updates and inspiration.

\medskip
\topic{Financial Overview and Significant Changes}
\noindent
The Treasurer’s report confirmed a healthy financial year, continuing the club’s steady growth in membership and income. Total income for 2025 reached £6387.21, with total expenses at £4881.45, giving a net surplus of £1505.76. Club assets now stand at just over £7200, including cash, bank holdings, and tangible items such as tables and trees.

The AGM introduced several important financial changes for the coming membership year. The most significant is the shift of membership renewal from January to April. This realignment reduces administrative pressure, places renewals away from the holiday season, and creates a smoother link with the club’s financial year. All existing memberships will therefore be extended until April 2026.

Membership fees will increase from £30 to £40, largely due to rising hall hire costs. To counterbalance this, members paying upfront will now have all guest speaker top-up fees included in their annual fee. The “12 for 10” model for club nights remains in place, offering good value for frequent attendees. Other changes include renaming joint membership to “Family membership” (with a £10 reduction) and offering 10% discounts for full-time students. A formal mid-year reduction point has also been set: from October onward, new members will pay the reduced £20 rate.

These adjustments aim to simplify administration while providing greater value back to members. The Treasurer also outlined projected increases in financial stability for 2025 and 2026, anticipating small surplus levels that allow the club to absorb costs and plan future improvements.

\medskip
\topic{Activities, Feedback, and Member Experience}
\noindent
Feedback gathered during the AGM reflected widespread satisfaction with the club’s activities through the year. Members consistently highlighted the warm atmosphere, the breadth of shared knowledge, and the variety of events. The committee encourages anyone with suggestions or concerns to speak with them at any time; maintaining open communication remains a priority.
%
For members not using WhatsApp or Discord, communication by email will continue, ensuring no one is left out of important announcements or club business.

\medskip
\topic{Allotment, Shows, and Subcommittees}
\noindent
Several areas of club activity saw updates. The allotment, overseen by Brendan and Steve, now has additional space after the clearing of older material. The intention is that members growing trees there plant a second tree on behalf of the club, with labelled areas and a small charge to contribute to site upkeep. A future auction of allotment material is planned to help cover running costs.
%
Preparations for next year's shows were also discussed. The display team remains in place, and tree submissions for the club show will, where relevant, also serve as submissions for the prestigious Swindon Show. The Swindon display will consist of two six-foot tables, likely accommodating five trees. Members submitting work must be ready to assist with setup and takedown.
%
The club will again include suiseki and shohin displays, with Paul and Rory proposed to lead their respective areas. A backdrop for the shohin section will be sourced ahead of time. Meanwhile, the sales table, refreshments, and progression table teams will continue in their established roles.

\medskip
\topic{New Business and Member Opportunities}
\noindent
Interest was raised in revisiting the outdoor display at Strawberry Hill House, and the committee will investigate possible dates. Brendan also requested feedback on Discord structure, noting that the current number of channels may warrant simplification.
%
The New Talent competition was again highlighted, with members encouraged to consider applying. The newsletter also received strong praise, but contributions from members are needed to maintain its quality. Even brief notes or outlines from events are valuable, and the editor will happily transform them into polished pieces.

\medskip
\topic{Closing Remarks}
\noindent
The meeting closed with a reaffirmation of members’ enthusiasm and confidence in the club’s journey. With stable leadership, thoughtful financial restructuring, and active participation across committees and events, the club enters the year ahead in excellent shape and with a spirit of shared purpose.

\closearticle
